\documentclass[11pt, a4paper]{article}

\usepackage{amsmath, amssymb}
\usepackage{graphicx}
\usepackage{booktabs}
\usepackage{hyperref}
\usepackage{geometry}
\usepackage{parskip}
\usepackage{caption}
\usepackage{siunitx}
\usepackage{xcolor}

\geometry{margin=2.5cm}

\title{Radial Current Distribution in a Magnetohydrodynamic Corbino Disk:\\
       Fall-Off Length Scaling with Applied Field and Fluid Height}
\author{PHYS 492 Project}
\date{\today}

\begin{document}

\maketitle

\begin{abstract}
We post-process COMSOL finite-element simulations of liquid gallium in a Corbino disk
driven by a radial current under a uniform axial magnetic field.
Two independent parametric sweeps — over applied field $B_0 \in [0.5,\,14]$\,T at
fixed fluid height $h = \SI{1}{mm}$, and over $h \in [0.5,\,15]$\,mm at fixed
$B_0 = \SI{5}{T}$ — both characterise the same quantity: the radial length $\delta$
over which $J_r$ transitions from its near-electrode value to the bulk.
From the B-sweep we find $\delta \propto B_0^{-0.43\pm0.01}$; from the H-sweep,
$\delta \propto h^{0.52\pm0.01}$.
Both results are consistent with the single universal law
\[
   \delta \;=\; \alpha\,\sqrt{D(B_0)\,h}, \qquad \alpha = 1.41 \pm 0.02,
\]
where $D = \sqrt{\mu/(\sigma B_0^2)}$ is the Hartmann layer thickness.
The collapse of both datasets onto this line (RMS scatter 4.7\%) confirms that
$\sqrt{Dh}$ — the geometric mean of the two intrinsic length scales — governs the
current redistribution near the inner electrode.
\end{abstract}

% ─────────────────────────────────────────────────────────────────────────────
\section{Physical Setup}

The Corbino disk consists of liquid gallium filling an annular region between an inner
electrode at $r_\mathrm{in} = \SI{0.1}{mm}$ and an outer electrode at
$r_\mathrm{out} = \SI{4.5}{mm}$, with fluid height $h$.
A uniform axial magnetic field $B_0\hat{z}$ drives azimuthal rotation through the
body force $F_\phi = -J_r B_0$.
No-slip boundary conditions at the top and bottom walls ($z = 0$ and $z = h$) create
Hartmann boundary layers of thickness
\begin{equation}
    D(B_0) = \sqrt{\frac{\mu}{\sigma B_0^2}},
    \label{eq:D}
\end{equation}
where $\mu = \SI{1.95e-3}{Pa.s}$ and $\sigma = \SI{3.7e6}{S/m}$ for liquid gallium.
The Hartmann number $\mathrm{Ha} = B_0 h\sqrt{\sigma/\mu} = h/D$ grows with both
$B_0$ and $h$; across the two sweeps it spans $\mathrm{Ha} \in [109,\,3267]$.

At large $\mathrm{Ha}$, the net radial current concentrates in the thin Hartmann
layers near $z=0$ and $z=h$.
The midplane ($z=h/2$) therefore carries a \emph{reverse} (negative) radial current
in the bulk — the return flow — while the net radial transport is outward.
This structure is visible throughout the H-sweep data and was observed for
$B_0 \gtrsim \SI{2}{T}$ in the B-sweep.

The COMSOL model is 2D axisymmetric, coupling Electric Currents and Laminar Flow
(swirl) with the Lorentz force $-J_r B_0$ and back-EMF
$\sigma U_\phi B_0$ applied manually.

% ─────────────────────────────────────────────────────────────────────────────
\section{Data}

\paragraph{B-sweep.}
$B_0 \in [0.5,\,14]$\,T (step \SI{0.5}{T}), 28 values, $h = \SI{1}{mm}$ fixed.
The mesh refines with $B_0$ (first element $\approx D/5$); row count grows from
685 (at \SI{0.5}{T}) to 2011 (at \SI{14}{T}).

\paragraph{H-sweep.}
$h \in [0.5,\,15]$\,mm (step \SI{0.5}{mm}), 30 values, $B_0 = \SI{5}{T}$ fixed.
This was confirmed by cross-referencing $J_r \cdot r$ at $r_\mathrm{in}$: the value
$2.194 \times 10^{-3}$\,A/m from the B-sweep at $B_0 = \SI{5}{T}$ agrees with the
H-sweep at $h = \SI{1}{mm}$ to 0.09\%, identifying $B_0$ without ambiguity.
The mesh has 949 rows for all $h$ values (designed to resolve Hartmann layers at
$B_0 = \SI{5}{T}$); the Hartmann number grows from $\mathrm{Ha}=109$ ($h=\SI{0.5}{mm}$)
to $\mathrm{Ha}=3267$ ($h=\SI{15}{mm}$).

Both exports are at the domain midplane $z=h/2$ along $r\in[r_\mathrm{in},\,\SI{3}{mm}]$.

% ─────────────────────────────────────────────────────────────────────────────
\section{Analysis Method}

The analysis follows the same procedure for both sweeps.

\paragraph{Bulk constant.}
At large $\mathrm{Ha}$ the midplane radial current in the bulk is
Hartmann-screened to near zero; $C_\mathrm{bulk} \approx 0$ throughout the H-sweep
($\mathrm{Ha} \geq 109$) and for $B_0 \geq \SI{2}{T}$ in the B-sweep.
The excess is therefore $\varepsilon(r) = J_r(r)\,r$ (B-sweep: minus a small residual
bulk constant estimated from the median of $J_r\cdot r$ at $r > 3r_\mathrm{in}$;
H-sweep: $C_\mathrm{bulk}\equiv 0$).

\paragraph{Decay and zero-crossing.}
For the H-sweep, the midplane $J_r$ reverses sign at an intermediate radius (the
Hartmann return current, Fig.~\ref{fig:jr_h}).
The decay search is therefore restricted to $r < r_\mathrm{zero}$, the first zero
crossing of $J_r\cdot r$ after the near-electrode peak, to avoid the outer electrode's
contribution.

\paragraph{Fall-off length.}
For each parameter slice we locate the peak excess $\varepsilon_\mathrm{peak}$
(in $r < 3r_\mathrm{in}$), interpolate onto a 500-point grid, and define
$\delta = r_{1/2} - r_\mathrm{in}$, where $r_{1/2}$ is the first $r$ at which
$\varepsilon$ drops to $\tfrac{1}{2}\varepsilon_\mathrm{peak}$.

% ─────────────────────────────────────────────────────────────────────────────
\section{Results}

\subsection{Current profiles}

Figure~\ref{fig:jr_b} shows $J_r(r)$ for five representative $B_0$ values from the
B-sweep.  The $1/r$ bulk form holds at low $B_0$; at higher $B_0$ the midplane current
collapses toward zero in the bulk.

Figure~\ref{fig:jr_h} shows $J_r(r)$ for five representative $h$ values from the
H-sweep.  The near-electrode excess decays, crosses zero, reaches a negative plateau
(Hartmann return current), and recovers near $r_\mathrm{out}$.  Larger $h$ (larger
$\mathrm{Ha}$) produces a deeper negative minimum and the zero crossing shifts outward.

\begin{figure}[ht]
    \centering
    \includegraphics[width=0.82\textwidth]{../figures/jr_profiles.pdf}
    \caption{B-sweep: $J_r(r)$ at five $B_0$ values.
             Dashed: $J_r\propto 1/r$ reference.  As $B_0$ rises, the midplane
             bulk current is suppressed by Hartmann screening.}
    \label{fig:jr_b}
\end{figure}

\begin{figure}[ht]
    \centering
    \includegraphics[width=\textwidth]{../figures/jr_profiles_h.pdf}
    \caption{H-sweep ($B_0 = \SI{5}{T}$): $J_r(r)$ at five $h$ values.
             \textit{Left}: near-electrode region.
             \textit{Right}: full profile showing the negative Hartmann return
             current and outer-electrode recovery.
             Dashed horizontal: $J_r = 0$.}
    \label{fig:jr_h}
\end{figure}

\subsection{B-sweep scaling}

Figure~\ref{fig:falloff_b} shows the normalised excess profiles from the B-sweep.
Figure~\ref{fig:scaling_b} plots $\delta_{1/2}$ versus $B_0$.
A power-law fit gives
\begin{equation}
    \delta_{1/2}^{(B)} = (190 \pm 2)\,\mu\mathrm{m} \cdot B_0^{-0.43\pm0.01}.
    \label{eq:fit_B}
\end{equation}

\begin{figure}[ht]
    \centering
    \includegraphics[width=0.82\textwidth]{../figures/falloff_profiles.pdf}
    \caption{B-sweep: normalised excess versus $r - r_\mathrm{in}$ for all $B_0$.
             Tick marks show $r_{1/2}$.  Higher $B_0$ → sharper decay.}
    \label{fig:falloff_b}
\end{figure}

\begin{figure}[ht]
    \centering
    \includegraphics[width=\textwidth]{../figures/scaling_r_half.pdf}
    \caption{B-sweep scaling: $\delta_{1/2}$ vs $B_0$.
             \textit{Left}: linear axes with fitted curve (Eq.~\ref{eq:fit_B}) and
             $1.41\sqrt{Dh}$ prediction.
             \textit{Right}: log-log confirmation of the power law.}
    \label{fig:scaling_b}
\end{figure}

\subsection{H-sweep scaling}

Figure~\ref{fig:falloff_h} shows the normalised excess profiles from the H-sweep.
Figure~\ref{fig:scaling_h} plots $\delta_{1/2}$ versus $h$.
A power-law fit gives
\begin{equation}
    \delta_{1/2}^{(h)} = (3275\,\mu\mathrm{m\,m}^{-n}) \cdot h^{0.52\pm0.01}.
    \label{eq:fit_h}
\end{equation}
The exponent $0.52$ is indistinguishable from $1/2$ within the mesh quantisation
uncertainty.  The amplitude ratio $\alpha = \delta/\sqrt{Dh}$ converges from
$1.35$ at $h=\SI{0.5}{mm}$ to $1.44$ at $h=\SI{15}{mm}$, consistent with
an asymptotic value of $1.41$ as $\mathrm{Ha}\to\infty$.

\begin{figure}[ht]
    \centering
    \includegraphics[width=0.82\textwidth]{../figures/falloff_profiles_h.pdf}
    \caption{H-sweep: normalised excess versus $r - r_\mathrm{in}$ for all $h$
             ($B_0 = \SI{5}{T}$).  Tick marks show $r_{1/2}$.
             Larger $h$ → wider profile, consistent with $\delta\propto h^{1/2}$.}
    \label{fig:falloff_h}
\end{figure}

\begin{figure}[ht]
    \centering
    \includegraphics[width=\textwidth]{../figures/scaling_h.pdf}
    \caption{H-sweep scaling: $\delta_{1/2}$ vs $h$.
             \textit{Left}: linear axes with fitted curve (Eq.~\ref{eq:fit_h})
             and $1.41\sqrt{Dh}$ prediction.
             \textit{Right}: log-log confirmation.}
    \label{fig:scaling_h}
\end{figure}

\subsection{Universal collapse}

Figure~\ref{fig:collapse} is the central result.
Both datasets are plotted as $\delta_{1/2}$ versus $\sqrt{D(B_0)\,h}$.
The B-sweep provides points where $D$ varies and $h$ is fixed; the H-sweep provides
points where $h$ varies and $D$ is fixed.
The two sweeps overlap at the cross-validation point ($B_0 = \SI{5}{T}$, $h = \SI{1}{mm}$,
marked with a star), where they agree to 1.8\%.

A free power-law fit to the combined dataset gives
\begin{equation}
    \delta_{1/2} \;\propto\; \bigl(\sqrt{Dh}\bigr)^{1.01\pm0.01},
    \label{eq:combined}
\end{equation}
and a forced-slope fit ($\delta = \alpha\sqrt{Dh}$) yields
\begin{equation}
    \boxed{\delta_{1/2} \;=\; (1.41 \pm 0.02)\,\sqrt{D(B_0)\,h}
           \;=\; (1.41 \pm 0.02)\left(\frac{\mu h^2}{\sigma B_0^2}\right)^{1/4},}
    \label{eq:result}
\end{equation}
with RMS scatter 4.7\% across 58 data points spanning a factor of six in $\sqrt{Dh}$.

\begin{figure}[ht]
    \centering
    \includegraphics[width=\textwidth]{../figures/combined_collapse.pdf}
    \caption{Universal scaling collapse.  B-sweep circles (blue, colour = $B_0$)
             and H-sweep squares (orange, colour = $h$) both lie on the line
             $\delta = 1.41\sqrt{Dh}$ (black).
             Star: cross-validation point ($B_0=5$ T, $h=1$ mm, 1.8\% agreement
             between sweeps).
             \textit{Left}: linear axes.  \textit{Right}: log-log.}
    \label{fig:collapse}
\end{figure}

% ─────────────────────────────────────────────────────────────────────────────
\section{Discussion}

\paragraph{Why $\sqrt{Dh}$?}
The fall-off length $\delta$ couples two distinct length scales.
The Hartmann layer $D$ controls how quickly the current redistributes in $z$: thinner
$D$ means faster $z$-redistribution, which confines the 2D near-electrode structure
to a shorter radial extent.
The fluid height $h$ sets the radial redistribution length because radial diffusion
and the $z$-averaged Lorentz body force both scale with $h$.
The geometric mean $\sqrt{Dh}$ arises when these two effects balance, as in many
mixed-boundary-layer problems.

\paragraph{The return current.}
At large $\mathrm{Ha}$, the midplane $J_r$ is negative in the bulk.
This is the Hartmann return current: the net outward current is carried in the
$z$-boundary layers, while the bulk (midplane) closes the current loop inward.
The H-sweep makes this structure explicit across 30 values of Ha.
Importantly, the fall-off length analysis is not contaminated by the return current
because the decay search is terminated at the first $J_r\cdot r = 0$ crossing.

\paragraph{Convergence of $\alpha$.}
The amplitude $\alpha = \delta/\sqrt{Dh}$ increases slightly with Ha
(from $1.35$ at $\mathrm{Ha}=109$ to $1.44$ at $\mathrm{Ha}=3267$).
This weak Ha-dependence suggests a small sub-leading correction of order $Ha^{-1/2}$
that vanishes in the $\mathrm{Ha}\to\infty$ limit where the asymptotic value
$\alpha\approx 1.41$ is recovered.

\paragraph{Limitations.}
The H-sweep mesh was designed for $B_0 = \SI{5}{T}$ and is under-resolved for large
$h$ (where $\mathrm{Ha}\gg 1$), potentially introducing mesh-quantisation artefacts
at high $h$ comparable to those observed in the B-sweep.
An $r_\mathrm{out}$-limited domain ($\SI{3}{mm}$ vs $\SI{4.5}{mm}$ nominal) means
the outer-electrode contribution enters the data at large $r$; the zero-crossing
cutoff mitigates this but may slightly underestimate $r_{1/2}$ at large $h$.

% ─────────────────────────────────────────────────────────────────────────────
\section{Summary}

\begin{itemize}
    \item Two independent parametric sweeps — varying $B_0$ at fixed $h$ and
          varying $h$ at fixed $B_0$ — both yield $\delta\propto (D h)^{1/2}$.
    \item The combined free-fit exponent on $\sqrt{Dh}$ is $1.01\pm0.01$,
          confirming linearity.
    \item The universal amplitude is $\alpha = 1.41\pm0.02$ across 58 data points
          spanning $\sqrt{Dh}\in[40,\,262]$\,µm.
    \item The midplane $J_r$ profile shows a clear Hartmann return current at
          intermediate $r$, with the negative minimum deepening with $\mathrm{Ha}$.
    \item The fall-off length $\delta = 1.41({\mu h^2}/{\sigma B_0^2})^{1/4}$
          is the governing length scale for near-electrode current redistribution
          in this MHD Corbino geometry.
\end{itemize}

\end{document}
